\begin{derivation}[电弱三规范玻色子顶角]
一般 Yang--Mills 三点张量由场强的非线性项固定。对 $SU(2)_L$，生成元满足

\[
 [t^a,t^b]=\ii\epsilon^{abc}t^c,
\]
所以结构常数就是 $f^{abc}=\epsilon^{abc}$。由于 $\epsilon^{abc}$ 在任意两个指标相同时为零，含一对 $W^+,W^-$ 的非零三点项必须由原始 $\mathcal W^1,\mathcal W^2$ 分量组成，而第三条原始中性规范场是 $\mathcal W^3$。把带电场定义的逆变换
\begin{equation*}
 \mathcal W^1=\frac{\mathcal W^++\mathcal W^-}{\sqrt2},
 \qquad
 \mathcal W^2=\frac{\ii(\mathcal W^+-\mathcal W^-)}{\sqrt2}
\end{equation*}
代入群因子。固定第一条外腿为 $W^+$、第二条为 $W^-$、第三条为 $\mathcal W^3$ 时，两个非零实基分量给
\begin{align*}
 \sum_{a,b=1}^{2}C_a^+C_b^-\epsilon^{ab3}
 &=\frac1{\sqrt2}\left(-\frac{\ii}{\sqrt2}\right)\epsilon^{123}
 +\frac{\ii}{\sqrt2}\frac1{\sqrt2}\epsilon^{213}\\
 &=-\frac{\ii}{2}-\frac{\ii}{2}\\
 &=-\ii.
\end{align*}
一般实基顶角的群因子是 $-g\epsilon^{ab3}$，所以带电基中 $W^+W^-\mathcal W^3$ 的总体因子变为 $+\ii g$。洛伦兹张量仍保持
\begin{equation*}
 \mathcal V_{\mu\nu\rho}(p,q,r)
 =\eta_{\mu\nu}(p-q)_\rho
 +\eta_{\nu\rho}(q-r)_\mu
 +\eta_{\rho\mu}(r-p)_\nu.
\end{equation*}
再用
\begin{equation*}
 \mathcal W^3=s_WA+c_W\mathcal Z
\end{equation*}
得到
\begin{equation*}
 g\mathcal W^3=gs_WA+gc_W\mathcal Z.
\end{equation*}
由\eqnrefs{eq:gem-definition-dp11} 的 $gs_W=g_{\rm em}$ 代入上式，得到
\begin{equation*}
 g_{WW\gamma}=g_{\rm em},
 \qquad
 g_{WWZ}=gc_W,
\end{equation*}
即\eqnrefs{eq:sm-wwv-couplings-dp11}。把同一 $\mathcal V_{\mu\nu\rho}$ 的三条动量依次指定为 $p_+,p_-,k$，并保持全部流入约定，便得到\eqnrefs{eq:sm-wwv-vertex-dp11}；其中显式的 $+\ii$ 正来自上述实基到带电基的复线性变换。
\end{derivation}

\begin{derivation}[带电矢量顶角的 Ward 恒等式]
正文\eqnrefs{eq:sm-wwv-vertex-dp11} 已固定三矢量顶角的动量结构。取 $V=\gamma$ 并用光子四动量收缩，再借助三条外腿全部流入的动量守恒，就能把结果整理成两条 $W$ 逆传播子之差，即正文\eqnrefs{eq:sm-wwgamma-ward-dp41}。

从\eqnrefs{eq:sm-wwv-vertex-dp11} 取 $V=\gamma$ 并与光子动量 $k^\rho$ 收缩：
\begin{align*}
 \frac{k^\rho\Gamma^{WW\gamma}_{\mu\nu\rho}}{+\ii g_{\rm em}}
 ={}&\eta_{\mu\nu}k\cdot(p_+-p_-)
 +k_\nu(p_--k)_\mu
 +k_\mu(k-p_+)_\nu.
\end{align*}
动量守恒 $k=-p_+-p_-$ 给出第一项
\begin{align*}
 k\cdot(p_+-p_-)
 &=(-p_+-p_-)\cdot(p_+-p_-)\\
 &=-p_+^2+p_-^2.
\end{align*}
后两项分别展开为
\begin{align*}
 k_\nu(p_--k)_\mu
 &=-(p_{+\nu}+p_{-\nu})(p_{+\mu}+2p_{-\mu}),\\
 k_\mu(k-p_+)_\nu
 &=(p_{+\mu}+p_{-\mu})(2p_{+\nu}+p_{-\nu}).
\end{align*}
把乘积逐项展开，$p_{+\mu}p_{-\nu}$ 与 $p_{-\mu}p_{+\nu}$ 的交叉项互相抵消，剩下
\begin{equation*}
 k_\nu(p_--k)_\mu+k_\mu(k-p_+)_\nu
 =p_{+\mu}p_{+\nu}-p_{-\mu}p_{-\nu}.
\end{equation*}
因此
\begin{align*}
 k^\rho\Gamma^{WW\gamma}_{\mu\nu\rho}
 =+\ii g_{\rm em}\Big[
 &(p_-^2\eta_{\mu\nu}-p_{-\mu}p_{-\nu})\\
 &-(p_+^2\eta_{\mu\nu}-p_{+\mu}p_{+\nu})
 \Big].
\end{align*}
在两个括号中同时减去相同的质量项 $m_W^2c^2\eta_{\mu\nu}$ 不改变差值，于是恰好得到
\begin{equation*}
 k^\rho\Gamma^{WW\gamma}_{\mu\nu\rho}
 =+\ii g_{\rm em}
 \left[\Pi_{\mu\nu}(p_-)-\Pi_{\mu\nu}(p_+)\right],
\end{equation*}
即\eqnrefs{eq:sm-wwgamma-ward-dp41}。若两条 $W$ 外腿在壳，则
\begin{equation*}
 p_\pm^2=m_W^2c^2,
 \qquad
 p_\pm\cdot\varepsilon_\pm=0,
\end{equation*}
所以 $\varepsilon_+^\mu\Pi_{\mu\nu}(p_+)=0$ 且 $\Pi_{\mu\nu}(p_-)\varepsilon_-^\nu=0$。这给出外部光子纵向替换后的零振幅，并同时检查了动量号、质量项与偏振条件。
\end{derivation}

\begin{derivation}[电弱真空附近的二次涨落与混合项]
规范固定要消去规范场与 Goldstone 场的二次混合项，因此先把协变导数在真空附近按场次数展开，保留决定自由传播子的二次项。正文采用
\begin{equation*}
\Phi
=\sqrt{\hbar c}
\begin{pmatrix}
\widehat\phi^+\\[2pt]
\bigl(v+\widehat h+\ii\widehat\chi\bigr)/\sqrt2
\end{pmatrix}.
\end{equation*}
只保留 \(D_\mu\Phi\) 中对涨落场至多一次的部分。带电上分量可写成
\begin{equation*}
(D_\mu\Phi)_1^{(1)}
=\sqrt{\hbar c}\left(
\partial_\mu\widehat\phi^+
+\ii\kappa_W\mathcal W_\mu^+
\right),
\end{equation*}
其中 \(\kappa_W=m_Wc/\hbar\) 按正文定义。其模平方给
\begin{align*}
\frac{|(D_\mu\Phi)_1^{(1)}|^2}{\hbar c}
={}&\partial_\mu\widehat\phi^-\partial^\mu\widehat\phi^+
+\kappa_W^2\mathcal W^-_\mu\mathcal W^{+\mu}\\
&+\ii\kappa_W
\left(
\mathcal W^-_\mu\partial^\mu\widehat\phi^+
-\mathcal W^+_\mu\partial^\mu\widehat\phi^-
\right).
\end{align*}
因此带电 Goldstone 的动能、\(W\) 的质量二次型以及 \(W\)--Goldstone 导数混合来自同一个平方，三者并不是三个独立假设。

中性下分量的一阶部分同理为
\begin{equation*}
(D_\mu\Phi)_2^{(1)}
=\sqrt{\frac{\hbar c}{2}}
\left[
\partial_\mu\widehat h
+\ii\partial_\mu\widehat\chi
+\ii\kappa_Z\mathcal Z_\mu
\right],
\end{equation*}
其中真空方向对光子组合的系数为零，因此 \(A_\mu^{\rm geo}\) 不产生质量项。取模平方并只保留二次项，实部与虚部分离后得到
\begin{equation*}
\frac12(\partial\widehat h)^2
+\frac12(\partial\widehat\chi)^2
+\frac12\kappa_Z^2\mathcal Z_\mu\mathcal Z^\mu
+\kappa_Z\mathcal Z_\mu\partial^\mu\widehat\chi.
\end{equation*}
这一步同时说明为什么 \(Z\)--Goldstone 混合只涉及 \(\widehat\chi\)，而径向模 \(\widehat h\) 不与 \(Z_\mu\) 发生双线性导数混合。

最后处理势能。若真空由 \(V'(v)=0\) 固定，则在径向方向 Taylor 展开
\begin{equation*}
V(v+\widehat h)
=V(v)+\frac12V''(v)\widehat h^2
+\mathcal O(\widehat h^3),
\end{equation*}
一次项因驻点条件消失。用正文定义 \(\kappa_h^2\equiv V''(v)/(\hbar c)\) 后，拉氏量中的势能贡献为
\begin{equation*}
-\frac12\kappa_h^2\widehat h^2.
\end{equation*}
把带电、中性与径向三部分相加，就逐项得到正文\eqnrefs{eq:sm-rxi-quadratic-before-gf-dp11}。该二次核明确暴露出随后选择 \(R_\xi\) 规范固定函数所要消去的两类导数混合项。
\end{derivation}

\begin{derivation}[规范固定后的 Goldstone 质量]
先展开带电规范固定项的交叉部分：

\begin{align*}
 -\frac1{\xi_W}F^+F^-
 \supset{}&-\ii\kappa_W\widehat\phi^-\partial\!\cdot\mathcal W^+
 +\ii\kappa_W\widehat\phi^+\partial\!\cdot\mathcal W^-.
\end{align*}
对作用量积分分部，并假设场在积分边界足够快衰减，使表面项为零：
\begin{align*}
 \int\dd^4 x\,
 \widehat\phi^-\partial_\mu\mathcal W^{+\mu}
 &=-\int\dd^4 x\,
 (\partial_\mu\widehat\phi^-)\mathcal W^{+\mu},\\
 \int\dd^4 x\,
 \widehat\phi^+\partial_\mu\mathcal W^{-\mu}
 &=-\int\dd^4 x\,
 (\partial_\mu\widehat\phi^+)\mathcal W^{-\mu}.
\end{align*}
因此规范固定项贡献
\begin{equation*}
 +\ii\kappa_W\mathcal W^+\!\cdot\partial\widehat\phi^-
 -\ii\kappa_W\mathcal W^-\!\cdot\partial\widehat\phi^+,
\end{equation*}
恰与\eqnrefs{eq:sm-rxi-quadratic-before-gf-dp11} 的带电混合项逐项相消。中性部分为
\begin{equation*}
 -\frac1{2\xi_Z}(F^Z)^2
 \supset
 +\kappa_Z\widehat\chi\,\partial\!\cdot\mathcal Z
 \longrightarrow
 -\kappa_Z\mathcal Z\!\cdot\partial\widehat\chi,
\end{equation*}
同样消去 $Z$--$\chi$ 混合。

规范固定函数平方还留下纯标量项
\begin{equation*}
 -\xi_W\kappa_W^2\widehat\phi^+\widehat\phi^-
 -\frac12\xi_Z\kappa_Z^2\widehat\chi^2.
\end{equation*}
与标量传播子的标准二次质量项比较，得到
\begin{equation*}
 m_{\phi^\pm}^2=\xi_Wm_W^2,
 \qquad
 m_\chi^2=\xi_Zm_Z^2,
\end{equation*}
即\eqnrefs{eq:sm-rxi-goldstone-masses-dp11}。
\end{derivation}

\begin{derivation}[规范轨道导数与鬼场质量]
由正文的 $\mathcal M_{\rm FP}^{ab}(x,y)=\delta F^a(x)/\delta\alpha^b(y)$ 出发。在线性化真空附近，破缺规范变换使 Goldstone 场沿规范轨道移动：

\begin{equation*}
 \delta\widehat\phi^\pm
 =\pm\ii\kappa_W\alpha^\pm+\order(\text{场}),
 \qquad
 \delta\widehat\chi
 =\kappa_Z\alpha^Z+\order(\text{场}).
\end{equation*}
同时规范场的线性变换给
\begin{equation*}
 \delta(\partial\!\cdot\mathcal W^\pm)=\partial^2\alpha^\pm+\order(\text{场}),
 \qquad
 \delta(\partial\!\cdot\mathcal Z)=\partial^2\alpha^Z+\order(\text{场}),
\end{equation*}
以及 $\delta(\partial\!\cdot A^{\rm geo})=\partial^2\alpha^\gamma+\order(\text{场})$。代入\eqnrefs{eq:sm-rxi-gauge-functions-dp11}：
\begin{align*}
 \delta F^\pm
 &=\left(\partial^2+\xi_W\kappa_W^2\right)\alpha^\pm
 +\order(\text{场}),\\
 \delta F^Z
 &=\left(\partial^2+\xi_Z\kappa_Z^2\right)\alpha^Z
 +\order(\text{场}),\\
 \delta F^\gamma
 &=\partial^2\alpha^\gamma+\order(\text{场}).
\end{align*}
于是鬼场二次算符分别是
\begin{equation*}
 \partial^2+\xi_W\kappa_W^2,
 \qquad
 \partial^2+\xi_Z\kappa_Z^2,
 \qquad
 \partial^2,
\end{equation*}
与自由标量二次核比较便得到\eqnrefs{eq:sm-rxi-ghost-masses-dp11}。
\end{derivation}

\begin{derivation}[横纵投影给出的有质量矢量传播子]
取物理四动量 $p^\mu$，并记

\begin{equation*}
 P^T_{\mu\nu}=\eta_{\mu\nu}-\frac{p_\mu p_\nu}{p^2},
 \qquad
 P^L_{\mu\nu}=\frac{p_\mu p_\nu}{p^2}.
\end{equation*}
它们满足
\begin{equation*}
 (P^T)^2=P^T,
 \qquad
 (P^L)^2=P^L,
 \qquad
 P^TP^L=0,
 \qquad
 P^T+P^L=I.
\end{equation*}
规范固定后的约化二次核可写成
\begin{equation*}
 K_{\mu\nu}(p)
 =-(p^2-m_V^2c^2)\eta_{\mu\nu}
 +\left(1-\frac1{\xi_V}\right)p_\mu p_\nu.
\end{equation*}
用 $\eta=P^T+P^L$ 与 $p_\mu p_\nu=p^2P^L_{\mu\nu}$ 分解：
\begin{align*}
 K_{\mu\nu}
 ={}&-(p^2-m_V^2c^2)P^T_{\mu\nu}\\
 &-\frac1{\xi_V}(p^2-\xi_Vm_V^2c^2)P^L_{\mu\nu}.
\end{align*}
两个投影子空间互相正交，所以逆核逐块取倒数：
\begin{equation*}
 K^{-1}_{\mu\nu}
 =-\frac{P^T_{\mu\nu}}{p^2-m_V^2c^2}
 -\frac{\xi_VP^L_{\mu\nu}}{p^2-\xi_Vm_V^2c^2}.
\end{equation*}
乘上 Feynman 边界条件对应的 $\ii$ 并重新合并两个投影算符，得到
\begin{equation*}
 D^{(V)}_{\mu\nu}(p)
 =\frac{-\ii}{p^2-m_V^2c^2+\ii0^+}
 \left[
 \eta_{\mu\nu}
 -(1-\xi_V)
 \frac{p_\mu p_\nu}{p^2-\xi_Vm_V^2c^2+\ii0^+}
 \right],
\end{equation*}
即\eqnrefs{eq:sm-rxi-massive-vector-propagator-dp41}。$\xi_V=1$ 时纵向附加项消失；$\xi_V\to\infty$ 时传播子化为幺正规范形式。这两个极限分别检验了纵向极点结构与质量项的归一化。
\end{derivation}

\begin{derivation}[Higgs 势的真空与质量]
把正文 Higgs 势写成单一规范不变量
\begin{equation*}
 x\equiv\Phi^\dagger\Phi\ge0,
 \qquad
 V(x)=-\mu_H^2x+\frac{\lambda_H}{\hbar c}x^2.
\end{equation*}
对复二重态作任意变分时，$\delta x=\delta\Phi^\dagger\Phi+\Phi^\dagger\delta\Phi$，因此驻点满足
\begin{equation*}
 \frac{\partial V}{\partial\Phi^\dagger}
 =\left(-\mu_H^2+\frac{2\lambda_H}{\hbar c}\Phi^\dagger\Phi\right)\Phi=0.
\end{equation*}
除原点 $\Phi=0$ 外，非零驻点组成
\begin{equation*}
 \Phi^\dagger\Phi
 =\frac{\mu_H^2\hbar c}{2\lambda_H}.
\end{equation*}
当 $\mu_H^2>0$、$\lambda_H>0$ 时，原点沿径向的二阶变化为负，而上述非零流形是稳定极小值。利用电弱规范变换可以从这条真空轨道中选取代表元
\begin{equation*}
 \langle\Phi\rangle
 =\frac1{\sqrt2}\binom0v,
\end{equation*}
于是
\begin{equation*}
 \frac{v^2}{2}
 =\frac{\mu_H^2\hbar c}{2\lambda_H},
 \qquad
 v^2=\frac{\mu_H^2\hbar c}{\lambda_H},
 \qquad
 v_E\equiv\sqrt{\hbar c}\,v,
\end{equation*}
即正文\eqnrefs{eq:sm-higgs-vev-dp11}。

真空附近先保留一般规范轨道方向，写成
\begin{equation*}
 \Phi(x)=\frac1{\sqrt2}U_\xi(x)\binom0{v+h(x)}.
\end{equation*}
势能只依赖 $\Phi^\dagger\Phi$，所以 $U_\xi^\dagger U_\xi=I$ 使三个切向场 $\xi^I$ 完全不进入 $V$；势能只约束径向变量 $h$。取幺正规范 $U_\xi=I$ 后，
\begin{equation*}
 \Phi^\dagger\Phi=\frac12(v+h)^2,
\end{equation*}
代回势能：
\begin{align*}
 V(h)
 &=-\frac{\mu_H^2}{2}(v+h)^2
 +\frac{\lambda_H}{4\hbar c}(v+h)^4\\
 &=-\frac{\mu_H^2}{2}
 (v^2+2vh+h^2)
 +\frac{\lambda_H}{4\hbar c}
 (v^4+4v^3h+6v^2h^2+4vh^3+h^4).
\end{align*}
逐阶收集：
\begin{align*}
 V(h)=V_0
 &+v\left(-\mu_H^2+\frac{\lambda_Hv^2}{\hbar c}\right)h\\
 &+\left(-\frac{\mu_H^2}{2}+\frac{3\lambda_Hv^2}{2\hbar c}\right)h^2
 +\frac{\lambda_Hv}{\hbar c}h^3
 +\frac{\lambda_H}{4\hbar c}h^4.
\end{align*}
利用真空条件 $\lambda_Hv^2/(\hbar c)=\mu_H^2$，线性项消失，二次项化为
\begin{equation*}
 V(h)=V_0+\mu_H^2h^2
 +\frac{\lambda_Hv}{\hbar c}h^3
 +\frac{\lambda_H}{4\hbar c}h^4,
\end{equation*}
即正文\eqnrefs{eq:sm-higgs-potential-expanded-dp11}。

最后核对径向场的动能归一化。忽略规范场时，
\begin{equation*}
 (\partial_\mu\Phi_U)^\dagger(\partial^\mu\Phi_U)
 =\frac12\partial_\mu h\,\partial^\mu h,
\end{equation*}
所以 $h$ 已经是标准规范化的实标量。标准质量项为
\begin{equation*}
 -\frac12\left(\frac{m_hc}{\hbar}\right)^2h^2.
\end{equation*}
由于拉格朗日量含 $-V$，比较二次项得到
\begin{equation*}
 \frac12\left(\frac{m_hc}{\hbar}\right)^2=\mu_H^2.
\end{equation*}
再用
\begin{equation*}
 \mu_H^2=\frac{\lambda_Hv^2}{\hbar c}
 =\frac{\lambda_Hv_E^2}{(\hbar c)^2},
\end{equation*}
便有
\begin{equation*}
 (m_hc^2)^2=2\lambda_Hv_E^2,
 \qquad
 \boxed{m_hc^2=\sqrt{2\lambda_H}\,v_E},
\end{equation*}
即正文\eqnrefs{eq:sm-higgs-mass-dp11}。同一展开同时固定 $h^3$ 和 $h^4$ 自耦，因此质量、三点自耦和四点自耦都来自同一个径向势，而不是彼此独立的参数。
\end{derivation}
